\begin{obeactivity}
\textbf{Aktivitas 1.1: Enkripsi di penyimpanan dan saat dikirim}
Amati dua penggunaan enkripsi berikut.
\begin{enumerate}
    \item Enkripsi dokumen di dalam penyimpanan: buat satu berkas teks, lalu
          lindungi dengan aplikasi enkripsi berkas atau fitur kata sandi pada
          aplikasi pengolah dokumen.
    \item Enkripsi pesan yang dikirim: amati keterangan enkripsi ujung-ke-ujung
          pada aplikasi perpesanan yang Anda pakai.
\end{enumerate}
Tuliskan untuk masing-masing: apa yang dilindungi, siapa yang memegang kunci,
dan apa yang terjadi bila kunci hilang.
\end{obeactivity}

\begin{obeactivity}
\textbf{Aktivitas 1.2: Membedakan \textit{encoding}, \textit{hashing}, dan enkripsi}
Ambil sebuah kalimat pendek, lalu lakukan tiga hal terhadapnya dengan perkakas
bebas yang tersedia (misalnya perintah \code{base64}, fungsi hash bawaan sistem,
dan aplikasi enkripsi berkas).
\begin{enumerate}
    \item Catat panjang keluaran masing-masing. Apakah panjangnya bergantung pada
          panjang masukan?
    \item Coba balikkan tiap keluaran. Mana yang berhasil dikembalikan tanpa kunci?
    \item Ubah satu huruf pada kalimat asal, lalu ulangi ketiga proses. Bandingkan
          seberapa besar perubahannya.
\end{enumerate}
Simpulkan primitif mana yang cocok untuk kerahasiaan dan mana yang cocok untuk
integritas.
\end{obeactivity}

\begin{obeactivity}
\textbf{Aktivitas 1.3: Menelusuri sebuah insiden kebocoran data}
Pilih satu kasus kebocoran data di Indonesia (misalnya BPJS Kesehatan 2021,
Tokopedia 2020, atau BSI 2023). Telusuri pemberitaannya, lalu jawab pertanyaan
berikut secara tertulis.
\begin{enumerate}
    \item Berapa banyak data yang terdampak, dan jenis data apa saja yang bocor?
    \item Layanan kriptografi mana yang seharusnya melindungi data itu?
    \item Apa yang dapat dilakukan pengelola data untuk mencegah pengulangan?
\end{enumerate}
Presentasikan hasilnya secara singkat di depan kelas.
\end{obeactivity}
